\documentclass[11pt,a4paper]{article}
\usepackage[utf8]{inputenc}
\usepackage[british]{babel}
\usepackage[a4paper,margin=2.3cm]{geometry}
\usepackage{amsmath,amssymb}
\usepackage{graphicx}
\usepackage{tikz}
\usetikzlibrary{calc,angles,quotes,patterns}
\usepackage[most]{tcolorbox}
\usepackage{enumitem}
\usepackage{array}
\usepackage{titlesec}
\usepackage[hidelinks]{hyperref}

\definecolor{bandgreen}{RGB}{46,139,87}
\definecolor{bandblue}{RGB}{30,100,190}
\definecolor{bandamber}{RGB}{220,140,20}
\definecolor{bandpink}{RGB}{210,70,140}
\definecolor{bandpurple}{RGB}{110,50,160}

\newcounter{prob}
% \prob{label}{title}{source}
\newcommand{\prob}[3]{\par\bigskip\refstepcounter{prob}\label{#1}%
  \noindent\textbf{\theprob.\ #2}\hfill{\small\textit{#3}}\par\nopagebreak\smallskip}
\newcommand{\ans}[2]{\par\smallskip\noindent\textbf{\ref{#1}.}\ #2}

\newcommand{\band}[3]{%
  \clearpage
  \begin{tcolorbox}[colback=#1!12,colframe=#1,boxrule=1pt,arc=2pt]
  {\Large\bfseries\color{#1} #2}\hfill{\large #3}
  \end{tcolorbox}}

\newtcolorbox{diarynote}[1]{colback=black!4,colframe=black!45,boxrule=0.6pt,
  arc=2pt,fonttitle=\bfseries\small,title={#1},fontupper=\small}

\setlength{\parindent}{0pt}
\setlength{\parskip}{3pt}
\definecolor{bandgrey}{RGB}{85,85,85}
\begin{document}

\begin{titlepage}
\centering
\vspace*{3cm}
{\Huge\bfseries The Ladies' Diary\par}
\vspace{0.4cm}
{\LARGE The last years, 1835--1840\par}
\vspace{0.8cm}
{\Large Twenty problems from the final six issues\par}
\vspace{2cm}
{\large\itshape Virtue and taste, with female softness join'd,\\
All that subdues and captivates mankind,\\
In Britain's matchless fair resplendent shine.\par}
\vspace{0.3cm}
{\small The Diary's title page, 1835\par}
\vfill
{\large Working draft\par}
\end{titlepage}

\section*{A word before the problems}

Leybourn's collection stops in 1816. The Diary did not. It ran for another
twenty-four years, under Olinthus Gregory of the Royal Military Academy, until the
Stationers' Company merged it with \textit{The Gentleman's Diary} in 1841. None of those
later questions is in Leybourn.

This booklet is made from one bound volume of the last six issues, for 1835 to 1840, the
132nd to the 137th almanacs ``published of this kind''. The copy is Harvard College
Library's, digitised by HathiTrust, and came to Harvard from the estate of a George
Eastwood. A Mr Geo.\ Eastwood of Dewsbury is among the solvers in 1835; they may be the
same man.

Each issue printed the answers to the previous year's fifteen questions and set fifteen
new ones, followed by an appendix of ``curious and valuable mathematical papers''. So a
single issue never holds a question together with its answer. The questions in this
volume run from 1584 to 1673. The answers to the last fifteen, set in 1840, were printed
in 1841 in \textit{The Lady's and Gentleman's Diary}, which we do not have; for those, and
for any answer we could not read, the answers here are our own.

By now the solvers were schoolmasters, surveyors, cadets and assistants at the Nautical
Almanac Office, with W.~S.~B.~Woolhouse and William Rutherford of Woolwich answering
almost everything. They came from further afield, too. Two questions in 1839 were set by
``M.~Z.~Ameen Bey, a scientific Egyptian now studying in England'', who appears in 1840 as
``Colonel Ameen Bey, Turkish Artillery''. Professor Charles Gill of St Paul's College, New
York, sent answers and a question.

And there were still women. A ``Miss L.~L.'' set one question in 1834 and another in 1836
(problem~\ref{z-equator}), and her name appears among the solvers again and again in 1837.
A Mary Forster answered in 1837 too. In the enigmas and verse, women are everywhere:
Winifred Waverton of Lambeth, Mrs and Miss Richardson of Vauxhall, Helen Ogden of Oldham,
Frances Plumbe of Bromley, Jane Sherwill of Plympton and many more.

\subsection*{How to read the labels}

\textit{Q1608} is the Diary's own question number. \textit{After Q1669} means the question has
been retold, or its general data replaced by a worked case; the answer at the back is then
for the retelling. The bands are those of the first book in this series. These issues have
nothing young enough for Green, and their weight is in Purple and beyond.

\subsection*{How this edition was made}
This edition was prepared with the assistance of Claude, an AI model made by Anthropic.
Under the editor's direction, Claude worked from HathiTrust's machine-read text of the
Harvard copy. It drafted the retellings, solutions, notes and introduction, recomputed
every answer, and typeset the text. Selection, editorial decisions and final
responsibility for the text are the editor's.

%================ BLUE ================
\band{bandblue}{Blue}{ages 8--11}

\prob{z-diagonals}{Counting diagonals}{Q1629, from Gill's \textit{Mathematical Miscellany}}
How many diagonals can be drawn in a hexagon? In an octagon? Find a rule for a figure
with $n$ sides.

%================ AMBER ================
\band{bandamber}{Amber}{ages 11--14}

\prob{z-shilling}{The shilling on the chessboard}{Q1597, by W.~F.}
A shilling is tossed and lands flat on a chessboard. The side of each square is twice the
diameter of the shilling. What is the chance that it lies wholly inside one square,
touching no line?

\prob{z-isosceles}{Three to two}{Q1614, by Thomas Todd}
In an isosceles triangle, the two equal sides together are as long as the base and the
height together. Show that the base is to the height as 3 is to 2.

%================ PINK ================
\band{bandpink}{Pink}{ages 14--16}

\prob{z-ortho}{The orthocentre and the cotangent}{Q1599, by the Rev.~John Sampson}
In a triangle $ABC$ the three altitudes meet at $P$. Show that $PC = AB\cot C$.

\prob{z-area}{An area from the incircle}{Q1600, by Thomas Baker}
The incircle of a triangle touches side $AB$ and divides it into two pieces. Show that the
product of the two pieces and $\cot\tfrac12C$ equals the area of the triangle.

\prob{z-midpoints}{Midpoints and contact points}{Q1635, by W.~F.~Harrison}
On each side of a triangle, the distance between the midpoint of the side and the point
where the incircle touches it is half the difference of the other two sides. Prove it.

\prob{z-altitudes}{The area from three altitudes}{Q1605, by W.~Rutherford}
The three altitudes of a triangle are $p_1$, $p_2$, $p_3$. Find its area. Work it out for
altitudes 12, 15 and 20.

\prob{z-flag}{The flag and the tower}{after Q1669, by W.~Rutherford}
A flagstaff $a$ feet high stands on top of a tower $b$ feet high. Where on the level ground
must an observer stand, with his eye $h$ feet above the ground, so that the tower and the
flagstaff look equally tall? Work it out for $a = 30$, $b = 20$, $h = 5$.

\prob{z-arbelos}{Two semicircles in a semicircle}{Q1668, by W.~S.~B.~Woolhouse}
Draw a semicircle. With the two ends of its diameter as centres, draw two more
semicircles, on the same side, that touch each other. Show that a common tangent to these
two also touches the first semicircle.

\prob{z-equilateral}{Three equilateral triangles}{Q1633, by W.~Rutherford}
On each side of a triangle, draw an equilateral triangle outwards. From each corner of the
original triangle, drop a perpendicular to the line joining the two outer vertices
nearest to it. Show that the three perpendiculars meet in a point.

%================ PURPLE ================
\band{bandpurple}{Purple}{ages 16--17}

\prob{z-gauss}{Ameen Bey's integral}{Q1653, by M.~Z.~Ameen Bey}
Find $\displaystyle\int_0^\infty e^{-x^2}\,dx$.

\prob{z-series}{Ameen Bey's series}{Q1654, by M.~Z.~Ameen Bey}
Find the $n$th term, and the sum of $n$ terms, of $1 + 3 + 8 + 24 + 65 + 153 + 320 + 610 +
\cdots$.

\prob{z-cylinders}{Cylinders without end}{Q1663, by John Butterworth}
Cut the largest possible cylinder out of a cone. Then cut the largest cylinder out of the
small cone left on top, and so on for ever. What fraction of the cone do all the
cylinders fill?

\prob{z-billiard}{The shortest quadrilateral in a rectangle}{Q1621, by Robert Sumner}
Show that among all quadrilaterals with one corner on each side of a rectangle, the least
perimeter equals the sum of the two diagonals. Show that there are infinitely many such
quadrilaterals, all parallelograms with sides parallel to the diagonals.

\prob{z-parabola}{Circles in a parabola}{Q1622, by W.~F.~Harrison}
A string of circles is placed inside a parabola, along its axis, each touching the
parabola and the one before. The first touches it at the vertex and has its centre at the
focus. Find the radius of the $n$th circle.

\prob{z-locus}{The third circle}{Q1661, by Philip Beecroft}
Two circles touch each other and touch a straight line at the fixed points $A$ and $B$. A
third, small circle touches the line and both circles. As the first two circles change
size, keeping their points of contact $A$ and $B$, find the path of the centre of the
third.

\prob{z-pressure}{Equal pressures}{Q1641, by Robert Sumner}
A rectangle of length $a$ is lowered into water, tilted at any angle, with one short edge
just at the surface. Divide it, by lines parallel to that edge, into $n$ strips that bear
equal pressures.

\prob{z-equator}{To Constantinople by the equator}{Q1608, by Miss L.~L.}
A traveller undertakes to go from London to Constantinople and to touch the equator on the
way. Treating the Earth as a sphere, find where he should touch it and how long his
journey is, if it is to be as short as possible. Take London at latitude $51^\circ33'$,
Constantinople at $41^\circ1'$, and the difference of longitude $28^\circ58'$, as the
Diary's solver did.

%================ BEYOND ================
\band{bandgrey}{Beyond A-level}{series and motion}

\prob{z-test}{The integral test}{Q1671, by Amicus}
Let $\phi$ be positive and decreasing on $x \ge 1$. Show that $\int_1^\infty\phi(x)\,dx$ is
finite or infinite according as $\phi(1) + \phi(2) + \phi(3) + \cdots$ converges or
diverges.

\prob{z-hollow}{The hollowed cone}{Q1589, by James Hann}
A solid cone swings as a pendulum about a horizontal axis through its vertex. A cone with
the same base, but with its vertex at the centre of gravity of the first, is hollowed out
of it. How much does this change the time of a swing?

%================ ANSWERS ================
\clearpage
\section*{Answers}

\ans{z-diagonals}{9 in a hexagon and 20 in an octagon. From each corner you can draw a
diagonal to every corner except itself and its two neighbours, so $n - 3$ of them. That
counts each diagonal twice, so there are $n(n-3)/2$.}

\ans{z-shilling}{One in four. The shilling lies inside a square exactly when its centre
lands at least a radius from every side, that is, inside a smaller square whose side is
the square's side less one diameter. Here that is half the side, so a quarter of the
area. Mr Mason of Scoulton gave the general rule: with squares $n$ diameters across, the
chance is $(n-1)^2/n^2$.}

\ans{z-isosceles}{Let the equal sides be $s$, the base $b$ and the height $h$. Then
$2s = b + h$ and $s^2 = h^2 + \tfrac14b^2$. Squaring the first and comparing,
$b^2 + 2bh + h^2 = 4h^2 + b^2$, so $2b = 3h$.}

\ans{z-ortho}{Let $CF$ be the altitude from $C$ and $AD$ the one from $A$. Triangles $CPD$
and $ABD$ are similar, because both are right-angled and $\angle PCD = \angle BAD$. So
$PC/AB = CD/AD = \cot C$. This is Rutherford's proof; Thomas Hindle and W.~F.~Harrison
gave others. Miss L.~L. and Mary Forster are among the solvers listed.}

\ans{z-area}{With the usual notation the pieces of $AB$ are $s-a$ and $s-b$, and
$\cot\tfrac12C = (s-c)/r$. So the product is $(s-a)(s-b)(s-c)/r$. By Heron's formula
$(s-a)(s-b)(s-c) = \Delta^2/s$, and $\Delta = rs$, so the product is $\Delta$. W.~Hopps
proved it with the fact that if three arcs make up a semicircle, the sum of their tangents
equals their product.}

\ans{z-midpoints}{On side $BC$ the contact point is $s - b$ from $B$, and the midpoint is
$\tfrac12a$ from $B$. Their distance is $|s - b - \tfrac12a| = \tfrac12|c - b|$.}

\ans{z-altitudes}{Write $u_i = 1/p_i$ and $\sigma = \tfrac12(u_1+u_2+u_3)$. The sides are
$2\Delta u_i$, and Heron's formula then gives
$1/\Delta = 4\sqrt{\sigma(\sigma-u_1)(\sigma-u_2)(\sigma-u_3)}$. For altitudes 12, 15 and
20 the area is 150: the triangle has sides 25, 20 and 15. Several solvers pointed out that
the question was already in Emerson's \textit{Algebra}.}

\ans{z-flag}{The line of sight to the top of the tower must bisect the angle between the
lines of sight to its foot and to the top of the flag. By the angle-bisector theorem, the
distances from the eye to the foot of the tower and to the top of the flag are as $b$ to
$a$. That gives
\[ x^2 = \frac{a^2h^2 - b^2(a+b-h)^2}{b^2 - a^2} \]
for the distance $x$ from the tower. With $a = 30$, $b = 20$, $h = 5$ it is about 39.7 feet,
where each subtends $27^\circ53'$. With the eye on the ground ($h = 0$) it is
$b\sqrt{(a+b)/(a-b)}$, possible only if the flagstaff is taller than the tower. The answer
was printed in 1841, in a volume we do not have; this one is ours.}

\ans{z-arbelos}{Let the first semicircle have centre $O$ and radius $R$, and the others
centres $A$, $B$ (the ends of the diameter) and radii $p$, $q$ with $p + q = 2R$. A common
tangent is at distances $p$ and $q$ from $A$ and $B$. Its distance from the midpoint $O$ of
$AB$ is the average, $\tfrac12(p+q) = R$. So it touches the first semicircle too. (The
answer appeared in 1841; this is ours.)}

\ans{z-equilateral}{Let $A'$, $B'$, $C'$ be the outer vertices on $BC$, $CA$, $AB$. The circle
with centre $B'$ through $A$ also passes through $C$, and the circle with centre $C'$
through $A$ also passes through $B$. These two circles meet at $A$, so their radical axis
is the line through $A$ perpendicular to $B'C'$: the first of the three perpendiculars.
In the same way the other two perpendiculars are the radical axes of the other pairs among
the three circles centred at $A'$, $B'$, $C'$. The three radical axes of three circles
always meet in a point, the radical centre.}

\ans{z-gauss}{$\tfrac12\sqrt\pi$. Square the integral and write it as a double integral
over the quarter-plane, $\iint e^{-(x^2+y^2)}\,dx\,dy$. In polar co-ordinates this is
$\int_0^{\pi/2}\int_0^\infty e^{-r^2}r\,dr\,d\theta = \tfrac\pi4$.}

\ans{z-series}{The fifth differences are constant, so the $n$th term is a polynomial of
degree five:
\[ u_n = \tfrac1{60}\bigl(n^5 + 15n^3 - 90n^2 + 254n - 120\bigr). \]
The next term is 1,081. The sum of $n$ terms is
$\tfrac1{360}\,n\bigl(n^5 + 3n^4 + 25n^3 - 135n^2 + 514n - 48\bigr)$. Professor Gill of
New York found the same differences, 2, 3, 8, 6, 2, and summed by the calculus of finite
differences.}

\ans{z-cylinders}{$\tfrac6{13}$ of the cone. The largest cylinder in a cone fills
$\tfrac49$ of it and is cut at a third of the height. The cone left on top has a third of
the height and a third of the radius, so $\tfrac1{27}$ of the volume. The cylinders fill
$\tfrac49\bigl(1 + \tfrac1{27} + \tfrac1{27^2} + \cdots\bigr) = \tfrac49\cdot\tfrac{27}{26} =
\tfrac6{13}$. The first step is Barbara Sidway's question of 1714. (The printed answer
appeared in 1841; this one is ours.)}

\ans{z-billiard}{Reflect the rectangle in its sides, as a billiard player would. A
quadrilateral with one corner on each side unfolds into a path from a point to its image
after four reflections, which is a translation by twice a diagonal. So its perimeter is at
least that distance, twice the length of a diagonal, which is the sum of the two
diagonals. The bound is reached when the path is straight: the sides are then parallel to
the diagonals, and any starting point on a side works.}

\ans{z-parabola}{Let the parabola be $y^2 = 4kx$, with focus $(k, 0)$. The first circle has
radius $k$. A circle centred on the axis at $(c, 0)$ with $c > 2k$ touches the parabola
when its radius satisfies $r^2 = 4k(c - k)$. Writing $r = 2ks$, the condition that each
circle touches the next gives $(s_{n+1} - 1)^2 = (s_n + 1)^2$ from the third circle on, and
$s_2 = 1 + \sqrt2$. So for $n \ge 2$ the radius of the $n$th circle is
$2k\bigl(2n - 3 + \sqrt2\bigr)$: from the second circle on, the radii go up in equal steps
of $4k$, twice the semi-parameter.}

\ans{z-locus}{An arc of a parabola through $A$ and $B$, with its vertex above the middle of
$AB$, at an eighth of $AB$. Put $AB = 2d$. Circles of radii $r_1$, $r_2$ touching a line
and each other have $AB^2 = 4r_1r_2$, so $r_1r_2 = d^2$; the small circle has
$1/\sqrt{r_3} = 1/\sqrt{r_1} + 1/\sqrt{r_2}$. Its centre is at height $y = r_3$ and
distance $x = 2\sqrt{r_1r_3}$ from $A$, and eliminating the radii gives
$y = x(2d - x)/4d$. (The answer appeared in 1841; this one is ours.)}

\ans{z-pressure}{The pressure on a strip is proportional to its area times the depth of its
centre, so the pressure on the part from the surface edge down to a distance $x$ along the
rectangle goes as $x^2$, whatever the tilt. The strips bear equal pressures when the
dividing lines are at distances $a\sqrt{1/n}, a\sqrt{2/n}, \dots, a\sqrt{(n-1)/n}$ from the
top edge.}

\ans{z-equator}{The shortest route makes equal angles with the equator on the way down and
the way up, like light reflected in a mirror. So reflect Constantinople in the equator, to
latitude $41^\circ1'$ south, and go by the great circle from London to that image. With the
Diary's figures the journey is $95^\circ56'$ of a great circle, about 6,660 miles, and the
traveller touches the equator $17^\circ11'$ of longitude east of London. William Rutherford
used the equal-angle principle. W.~Clough's printed answer is $96^\circ0'45''$, touching at
$16^\circ55'$.}

\ans{z-test}{Because $\phi$ is decreasing, on each interval $[k, k+1]$ we have
$\phi(k+1) \le \phi(x) \le \phi(k)$. Integrating and adding,
$\phi(2) + \cdots + \phi(n) \le \int_1^n\phi \le \phi(1) + \cdots + \phi(n-1)$. So the
integral and the series are bounded or unbounded together. This is the Maclaurin--Cauchy
integral test. (Amicus's own answer appeared in 1841.)}

\ans{z-hollow}{Let the cone have height $h$ and base radius $R$. Hung by its vertex, the
solid cone swings like a simple pendulum of length $(4h^2 + R^2)/5h$. The hollowed cone
swings like one of length $3(27h^2 + 8R^2)/110h$. So the time of a swing is multiplied by
\[ \sqrt{\frac{3(27h^2 + 8R^2)}{22(4h^2 + R^2)}}. \]
For a cone as wide as it is tall ($R = h$) this is 0.977, so it swings about 2.3 per cent
faster. The two times are equal when $R^2 = \tfrac72h^2$, and a wider cone swings more
slowly once hollowed. W.~S.~B.~Woolhouse answered it in 1836, by moments of inertia.}

\begin{diarynote}{When the Diary was wrong}
\textbf{Problem~\ref{z-equator}.} On his own figures, W.~Clough's track of
$96^\circ0'45''$ is about $4'$ of arc, some five miles, longer than the shortest, and his
point of contact is about $17'$ too far west. The editor printed several answers that
``differ only in numerical result on account of different assumptions of longitude and
latitude''. That is true of the others, but not of this difference, which is an error in
the working.
\end{diarynote}

\begin{diarynote}{Where the Diary simply worked differently}
\textbf{Problem~\ref{z-altitudes}.} Not an error, but not new: solvers pointed out that
the question had already been answered in Emerson's \textit{Algebra}.

\textbf{Not in this booklet.} Question 1623 asked for a proof that a certain sum of
rectangles about the centre of gravity of a triangle is constant. The editor printed a
note that ``some oversight appears to have been committed in the enunciation of this
question'', and gave instead the true value in terms of the sides.
\end{diarynote}

\vfill
{\small Source: \textit{The Ladies' Diary} for 1835--1840, Harvard College Library copy,
digitised by HathiTrust (\url{https://hdl.handle.net/2027/hvd.hxvmnt}).}

\end{document}
